% =========================================================
% ppt/src/minggu-01.tex
% Slide pertemuan 1
% Bab III — Pengantar Assembly dan Sistem Bilangan
%
% Sub-CPMK 1.1 | CPMK01 | CPL02
% =========================================================

\documentclass[aspectratio=169,11pt]{beamer}
% =========================================================
% ppt/src/preamble-ppt.tex
% Tema, warna, paket, dan makro bersama seluruh deck slide.
%
% Dipakai oleh setiap berkas minggu-NN.tex dengan pola:
%   \documentclass[aspectratio=169,11pt]{beamer}
%   % =========================================================
% ppt/src/preamble-ppt.tex
% Tema, warna, paket, dan makro bersama seluruh deck slide.
%
% Dipakai oleh setiap berkas minggu-NN.tex dengan pola:
%   \documentclass[aspectratio=169,11pt]{beamer}
%   % =========================================================
% ppt/src/preamble-ppt.tex
% Tema, warna, paket, dan makro bersama seluruh deck slide.
%
% Dipakai oleh setiap berkas minggu-NN.tex dengan pola:
%   \documentclass[aspectratio=169,11pt]{beamer}
%   \input{preamble-ppt}
%
% Palet warna disalin dari buku_ajar/config/styles.tex agar
% identitas visual slide sama dengan buku ajar.
% =========================================================

% ---------------------------------------------------------
% 1. Paket
% ---------------------------------------------------------
\usepackage[T1]{fontenc}
\usepackage[utf8]{inputenc}
\usepackage[indonesian]{babel}
\usepackage{lmodern}
\usepackage{tikz}
\usetikzlibrary{arrows.meta, positioning, calc, shapes.geometric}
\usepackage{booktabs}
\usepackage{array}
\usepackage{amsmath}
\usepackage{amsfonts}
\usepackage{pifont}          % \ding untuk tanda checklist
\usepackage{listings}

% ---------------------------------------------------------
% 2. Palet warna OBE (identik dengan buku)
% ---------------------------------------------------------
\definecolor{obeBlue}{HTML}{1F3864}
\definecolor{obeBlueLight}{HTML}{DCE4F5}
\definecolor{obeTeal}{HTML}{0F6E6E}
\definecolor{obeTealLight}{HTML}{D9F0EE}
\definecolor{obeGreen}{HTML}{2E6B34}
\definecolor{obeGreenLight}{HTML}{DEEFDF}
\definecolor{obeAmber}{HTML}{8A5A00}
\definecolor{obeAmberLight}{HTML}{FDF1D6}
\definecolor{obeMaroon}{HTML}{7A2233}
\definecolor{obeMaroonLight}{HTML}{F7E1E5}
\definecolor{obeGray}{HTML}{4A4A4A}
\definecolor{obeGrayLight}{HTML}{F0F0F0}
\definecolor{codeBack}{HTML}{F7F7F4}
\definecolor{codeFrame}{HTML}{C9C9C0}

% ---------------------------------------------------------
% 3. Identitas mata kuliah (dipakai seluruh deck)
% ---------------------------------------------------------
\newcommand{\NamaMK}{Bahasa Pemrograman Tingkat Rendah}
\newcommand{\NamaMKShort}{Bhs.\ Pemrograman Tingkat Rendah}
\newcommand{\KodeMK}{MWCA010}
\newcommand{\Prodi}{Program Studi Teknik Informatika}
\newcommand{\Fakultas}{Fakultas Teknologi Informasi}
\newcommand{\PerguruanTinggi}{Universitas Bale Bandung}
\newcommand{\DosenPengampu}{Cecep Suwanda, S.Si., M.Kom.}
\newcommand{\SemesterTeks}{Semester Ganjil}

% \PertemuanTeks dan \JudulBab ditetapkan di masing-masing
% berkas minggu-NN.tex sebelum \begin{document}.
\providecommand{\PertemuanTeks}{Pertemuan}

% ---------------------------------------------------------
% 4. Tema beamer
% ---------------------------------------------------------
\usetheme{default}
\setbeamertemplate{navigation symbols}{}

\setbeamercolor{normal text}{fg=black!88}
\setbeamercolor{structure}{fg=obeBlue}
\setbeamercolor{alerted text}{fg=obeMaroon}

% --- judul slide: pita biru penuh lebar ---
\setbeamercolor{frametitle}{fg=white,bg=obeBlue}
\setbeamerfont{frametitle}{series=\bfseries,size=\large}
\setbeamertemplate{frametitle}{%
    \nointerlineskip%
    \begin{beamercolorbox}[%
        wd=\paperwidth,ht=2.9ex,dp=1.5ex,
        leftskip=0.4cm,rightskip=0.4cm]{frametitle}%
        \usebeamerfont{frametitle}\insertframetitle%
    \end{beamercolorbox}%
    \nointerlineskip%
    \hbox{\textcolor{obeTeal}{\rule{\paperwidth}{1.1pt}}}%
}

% --- kaki slide: identitas mata kuliah, pertemuan, nomor ---
\setbeamercolor{footline}{fg=white,bg=obeGray}
\setbeamerfont{footline}{size=\scriptsize,family=\sffamily}
\setbeamertemplate{footline}{%
    \leavevmode%
    \hbox{%
        \begin{beamercolorbox}[wd=0.42\paperwidth,ht=2.1ex,dp=1.1ex,
            leftskip=0.4cm]{footline}%
            \usebeamerfont{footline}%
            \KodeMK\ \textbullet\ \NamaMKShort%
        \end{beamercolorbox}%
        \begin{beamercolorbox}[wd=0.32\paperwidth,ht=2.1ex,dp=1.1ex,
            center]{footline}%
            \usebeamerfont{footline}\PertemuanTeks%
        \end{beamercolorbox}%
        \begin{beamercolorbox}[wd=0.26\paperwidth,ht=2.1ex,dp=1.1ex,
            rightskip=0.4cm]{footline}%
            \usebeamerfont{footline}%
            \insertframenumber\,/\,\inserttotalframenumber%
        \end{beamercolorbox}%
    }%
    \vskip0pt%
}

% --- daftar ---
\setbeamercolor{itemize item}{fg=obeTeal}
\setbeamercolor{itemize subitem}{fg=obeTeal}
\setbeamercolor{enumerate item}{fg=obeBlue}
\setbeamertemplate{itemize item}{\raise0.2ex\hbox{\tiny$\blacksquare$}}
\setbeamertemplate{itemize subitem}{\raise0.2ex\hbox{\scriptsize$\bullet$}}

% --- kotak (pengganti tcolorbox pada buku) ---
\setbeamercolor{block title}{bg=obeBlueLight,fg=obeBlue}
\setbeamercolor{block body}{bg=obeBlueLight!45,fg=black!88}
\setbeamercolor{block title alerted}{bg=obeMaroonLight,fg=obeMaroon}
\setbeamercolor{block body alerted}{bg=obeMaroonLight!45,fg=black!88}
\setbeamercolor{block title example}{bg=obeTealLight,fg=obeTeal}
\setbeamercolor{block body example}{bg=obeTealLight!45,fg=black!88}

% ---------------------------------------------------------
% 5. Makro notasi (sama persis dengan buku)
%
% Catatan: \note TIDAK didefinisikan di sini karena beamer
% sudah memakai \note untuk catatan pembicara. Untuk catatan
% di slide dipakai \catatan.
% ---------------------------------------------------------
\newcommand{\code}[1]{\texttt{#1}}
\newcommand{\asm}[1]{\texttt{#1}}
\newcommand{\reg}[1]{\texttt{#1}}
\newcommand{\flag}[1]{\texttt{#1}}
\newcommand{\dir}[1]{\texttt{#1}}
\newcommand{\file}[1]{\texttt{#1}}
\newcommand{\eng}[1]{\textit{#1}}

\newcommand{\hex}[1]{\texttt{#1h}}
\newcommand{\bin}[1]{\texttt{#1b}}
\newcommand{\dec}[1]{\texttt{#1}}

\newcommand{\important}[1]{\textbf{Penting:} #1}
\newcommand{\catatan}[1]{\textbf{Catatan:} #1}
\newcommand{\term}[1]{\textbf{#1}}

\newcommand{\CPL}[1]{\textbf{CPL#1}}
\newcommand{\CPMK}[1]{\textbf{CPMK#1}}
\newcommand{\SubCPMK}[1]{\textbf{Sub-CPMK #1}}

\newcommand{\checkbox}{\ding{113}\hspace{0.35em}}
\newcommand{\checked}{\ding{51}\hspace{0.35em}}
\newcommand{\crossed}{\ding{55}\hspace{0.35em}}

\newcommand{\flowarrow}{\ensuremath{\rightarrow}}

% ---------------------------------------------------------
% 6. Makro tata letak slide
% ---------------------------------------------------------

% Keterangan sumber materi (ketertelusuran ke buku)
\newcommand{\Sumber}[1]{%
    \vfill
    {\scriptsize\color{obeGray}Sumber: #1\par}%
}

% Frame judul deck: \JudulDeck{<judul bab>}{<pertemuan>}
\newcommand{\JudulDeck}[2]{%
\begin{frame}[plain]
    \vspace{0.25cm}
    {\color{obeBlue}\sffamily\bfseries\small \NamaMK\ (\KodeMK)\par}
    \vspace{0.12cm}
    {\color{obeGray}\sffamily\scriptsize
     \Prodi\ \textbullet\ \Fakultas\ \textbullet\ \PerguruanTinggi\par}
    \vfill
    {\color{obeBlue}\sffamily\LARGE\bfseries #1\par}
    \vspace{0.4cm}
    {\color{obeTeal}\sffamily\large #2\par}
    \vspace{0.45cm}
    {\color{obeTeal}\rule{\textwidth}{1pt}}
    \vfill
    {\color{obeGray}\sffamily\small Dosen Pengampu\par}
    \vspace{0.08cm}
    {\sffamily\bfseries \DosenPengampu\par}
    \vspace{0.2cm}
    {\color{obeGray}\sffamily\small \SemesterTeks\par}
    \vspace{0.25cm}
\end{frame}
}

% Frame pemisah bagian
\newcommand{\PemisahBagian}[1]{%
\begin{frame}[plain]
    \vfill
    \begin{center}
        {\color{obeTeal}\sffamily\small\bfseries BAGIAN\par}
        \vspace{0.35cm}
        {\color{obeBlue}\sffamily\Large\bfseries #1\par}
        \vspace{0.35cm}
        {\color{obeTeal}\rule{0.55\textwidth}{0.9pt}}
    \end{center}
    \vfill
\end{frame}
}

% ---------------------------------------------------------
% 7. Gaya penampil kode
% ---------------------------------------------------------
\input{listing-ppt}

%
% Palet warna disalin dari buku_ajar/config/styles.tex agar
% identitas visual slide sama dengan buku ajar.
% =========================================================

% ---------------------------------------------------------
% 1. Paket
% ---------------------------------------------------------
\usepackage[T1]{fontenc}
\usepackage[utf8]{inputenc}
\usepackage[indonesian]{babel}
\usepackage{lmodern}
\usepackage{tikz}
\usetikzlibrary{arrows.meta, positioning, calc, shapes.geometric}
\usepackage{booktabs}
\usepackage{array}
\usepackage{amsmath}
\usepackage{amsfonts}
\usepackage{pifont}          % \ding untuk tanda checklist
\usepackage{listings}

% ---------------------------------------------------------
% 2. Palet warna OBE (identik dengan buku)
% ---------------------------------------------------------
\definecolor{obeBlue}{HTML}{1F3864}
\definecolor{obeBlueLight}{HTML}{DCE4F5}
\definecolor{obeTeal}{HTML}{0F6E6E}
\definecolor{obeTealLight}{HTML}{D9F0EE}
\definecolor{obeGreen}{HTML}{2E6B34}
\definecolor{obeGreenLight}{HTML}{DEEFDF}
\definecolor{obeAmber}{HTML}{8A5A00}
\definecolor{obeAmberLight}{HTML}{FDF1D6}
\definecolor{obeMaroon}{HTML}{7A2233}
\definecolor{obeMaroonLight}{HTML}{F7E1E5}
\definecolor{obeGray}{HTML}{4A4A4A}
\definecolor{obeGrayLight}{HTML}{F0F0F0}
\definecolor{codeBack}{HTML}{F7F7F4}
\definecolor{codeFrame}{HTML}{C9C9C0}

% ---------------------------------------------------------
% 3. Identitas mata kuliah (dipakai seluruh deck)
% ---------------------------------------------------------
\newcommand{\NamaMK}{Bahasa Pemrograman Tingkat Rendah}
\newcommand{\NamaMKShort}{Bhs.\ Pemrograman Tingkat Rendah}
\newcommand{\KodeMK}{MWCA010}
\newcommand{\Prodi}{Program Studi Teknik Informatika}
\newcommand{\Fakultas}{Fakultas Teknologi Informasi}
\newcommand{\PerguruanTinggi}{Universitas Bale Bandung}
\newcommand{\DosenPengampu}{Cecep Suwanda, S.Si., M.Kom.}
\newcommand{\SemesterTeks}{Semester Ganjil}

% \PertemuanTeks dan \JudulBab ditetapkan di masing-masing
% berkas minggu-NN.tex sebelum \begin{document}.
\providecommand{\PertemuanTeks}{Pertemuan}

% ---------------------------------------------------------
% 4. Tema beamer
% ---------------------------------------------------------
\usetheme{default}
\setbeamertemplate{navigation symbols}{}

\setbeamercolor{normal text}{fg=black!88}
\setbeamercolor{structure}{fg=obeBlue}
\setbeamercolor{alerted text}{fg=obeMaroon}

% --- judul slide: pita biru penuh lebar ---
\setbeamercolor{frametitle}{fg=white,bg=obeBlue}
\setbeamerfont{frametitle}{series=\bfseries,size=\large}
\setbeamertemplate{frametitle}{%
    \nointerlineskip%
    \begin{beamercolorbox}[%
        wd=\paperwidth,ht=2.9ex,dp=1.5ex,
        leftskip=0.4cm,rightskip=0.4cm]{frametitle}%
        \usebeamerfont{frametitle}\insertframetitle%
    \end{beamercolorbox}%
    \nointerlineskip%
    \hbox{\textcolor{obeTeal}{\rule{\paperwidth}{1.1pt}}}%
}

% --- kaki slide: identitas mata kuliah, pertemuan, nomor ---
\setbeamercolor{footline}{fg=white,bg=obeGray}
\setbeamerfont{footline}{size=\scriptsize,family=\sffamily}
\setbeamertemplate{footline}{%
    \leavevmode%
    \hbox{%
        \begin{beamercolorbox}[wd=0.42\paperwidth,ht=2.1ex,dp=1.1ex,
            leftskip=0.4cm]{footline}%
            \usebeamerfont{footline}%
            \KodeMK\ \textbullet\ \NamaMKShort%
        \end{beamercolorbox}%
        \begin{beamercolorbox}[wd=0.32\paperwidth,ht=2.1ex,dp=1.1ex,
            center]{footline}%
            \usebeamerfont{footline}\PertemuanTeks%
        \end{beamercolorbox}%
        \begin{beamercolorbox}[wd=0.26\paperwidth,ht=2.1ex,dp=1.1ex,
            rightskip=0.4cm]{footline}%
            \usebeamerfont{footline}%
            \insertframenumber\,/\,\inserttotalframenumber%
        \end{beamercolorbox}%
    }%
    \vskip0pt%
}

% --- daftar ---
\setbeamercolor{itemize item}{fg=obeTeal}
\setbeamercolor{itemize subitem}{fg=obeTeal}
\setbeamercolor{enumerate item}{fg=obeBlue}
\setbeamertemplate{itemize item}{\raise0.2ex\hbox{\tiny$\blacksquare$}}
\setbeamertemplate{itemize subitem}{\raise0.2ex\hbox{\scriptsize$\bullet$}}

% --- kotak (pengganti tcolorbox pada buku) ---
\setbeamercolor{block title}{bg=obeBlueLight,fg=obeBlue}
\setbeamercolor{block body}{bg=obeBlueLight!45,fg=black!88}
\setbeamercolor{block title alerted}{bg=obeMaroonLight,fg=obeMaroon}
\setbeamercolor{block body alerted}{bg=obeMaroonLight!45,fg=black!88}
\setbeamercolor{block title example}{bg=obeTealLight,fg=obeTeal}
\setbeamercolor{block body example}{bg=obeTealLight!45,fg=black!88}

% ---------------------------------------------------------
% 5. Makro notasi (sama persis dengan buku)
%
% Catatan: \note TIDAK didefinisikan di sini karena beamer
% sudah memakai \note untuk catatan pembicara. Untuk catatan
% di slide dipakai \catatan.
% ---------------------------------------------------------
\newcommand{\code}[1]{\texttt{#1}}
\newcommand{\asm}[1]{\texttt{#1}}
\newcommand{\reg}[1]{\texttt{#1}}
\newcommand{\flag}[1]{\texttt{#1}}
\newcommand{\dir}[1]{\texttt{#1}}
\newcommand{\file}[1]{\texttt{#1}}
\newcommand{\eng}[1]{\textit{#1}}

\newcommand{\hex}[1]{\texttt{#1h}}
\newcommand{\bin}[1]{\texttt{#1b}}
\newcommand{\dec}[1]{\texttt{#1}}

\newcommand{\important}[1]{\textbf{Penting:} #1}
\newcommand{\catatan}[1]{\textbf{Catatan:} #1}
\newcommand{\term}[1]{\textbf{#1}}

\newcommand{\CPL}[1]{\textbf{CPL#1}}
\newcommand{\CPMK}[1]{\textbf{CPMK#1}}
\newcommand{\SubCPMK}[1]{\textbf{Sub-CPMK #1}}

\newcommand{\checkbox}{\ding{113}\hspace{0.35em}}
\newcommand{\checked}{\ding{51}\hspace{0.35em}}
\newcommand{\crossed}{\ding{55}\hspace{0.35em}}

\newcommand{\flowarrow}{\ensuremath{\rightarrow}}

% ---------------------------------------------------------
% 6. Makro tata letak slide
% ---------------------------------------------------------

% Keterangan sumber materi (ketertelusuran ke buku)
\newcommand{\Sumber}[1]{%
    \vfill
    {\scriptsize\color{obeGray}Sumber: #1\par}%
}

% Frame judul deck: \JudulDeck{<judul bab>}{<pertemuan>}
\newcommand{\JudulDeck}[2]{%
\begin{frame}[plain]
    \vspace{0.25cm}
    {\color{obeBlue}\sffamily\bfseries\small \NamaMK\ (\KodeMK)\par}
    \vspace{0.12cm}
    {\color{obeGray}\sffamily\scriptsize
     \Prodi\ \textbullet\ \Fakultas\ \textbullet\ \PerguruanTinggi\par}
    \vfill
    {\color{obeBlue}\sffamily\LARGE\bfseries #1\par}
    \vspace{0.4cm}
    {\color{obeTeal}\sffamily\large #2\par}
    \vspace{0.45cm}
    {\color{obeTeal}\rule{\textwidth}{1pt}}
    \vfill
    {\color{obeGray}\sffamily\small Dosen Pengampu\par}
    \vspace{0.08cm}
    {\sffamily\bfseries \DosenPengampu\par}
    \vspace{0.2cm}
    {\color{obeGray}\sffamily\small \SemesterTeks\par}
    \vspace{0.25cm}
\end{frame}
}

% Frame pemisah bagian
\newcommand{\PemisahBagian}[1]{%
\begin{frame}[plain]
    \vfill
    \begin{center}
        {\color{obeTeal}\sffamily\small\bfseries BAGIAN\par}
        \vspace{0.35cm}
        {\color{obeBlue}\sffamily\Large\bfseries #1\par}
        \vspace{0.35cm}
        {\color{obeTeal}\rule{0.55\textwidth}{0.9pt}}
    \end{center}
    \vfill
\end{frame}
}

% ---------------------------------------------------------
% 7. Gaya penampil kode
% ---------------------------------------------------------
% =========================================================
% ppt/src/listing-ppt.tex
% Gaya penampil kode untuk slide.
%
% Diturunkan dari buku_ajar/config/listings.tex, tetapi
% disesuaikan untuk slide: ukuran lebih kecil (karena frame
% 16:9 dan kode harus utuh terlihat) dan jarak lebih rapat.
% =========================================================

% ---------------------------------------------------------
% 1. Definisi bahasa MASM/TASM 8086 (salinan dari buku)
% ---------------------------------------------------------
\lstdefinelanguage{MASM}{
    morekeywords={
        MOV, XCHG, LEA, LDS, LES, PUSH, POP, PUSHA, POPA, PUSHF, POPF,
        LAHF, SAHF, XLAT, IN, OUT, CBW, CWD,
        ADD, ADC, SUB, SBB, INC, DEC, NEG, CMP, MUL, IMUL, DIV, IDIV,
        AAA, AAS, AAM, AAD, DAA, DAS,
        AND, OR, XOR, NOT, TEST, SHL, SHR, SAL, SAR, ROL, ROR, RCL, RCR,
        JMP, JE, JZ, JNE, JNZ, JG, JGE, JL, JLE, JA, JAE, JB, JBE,
        JC, JNC, JS, JNS, JO, JNO, JP, JPE, JNP, JPO, JCXZ,
        LOOP, LOOPE, LOOPZ, LOOPNE, LOOPNZ, CALL, RET, INT, IRET, INTO,
        HLT, NOP, CLI, STI, CLD, STD, CLC, STC, CMC, WAIT, ESC, LOCK,
        DB, DW, DD, DQ, DT, EQU, ORG, SEGMENT, ENDS, ASSUME, END,
        PROC, ENDP, MACRO, ENDM, LOCAL, PUBLIC, EXTRN, EXTERN,
        INCLUDE, INCLUDELIB, MODEL, STACK, CODE, DATA, CONST, ALIGN, EVEN,
        OFFSET, PTR, BYTE, WORD, DWORD, QWORD, TBYTE,
        NEAR, FAR, SHORT, LABEL, TYPE, SIZE, LENGTH, DUP, THIS,
        AX, BX, CX, DX, AL, AH, BL, BH, CL, CH, DL, DH,
        SP, BP, SI, DI, IP, CS, DS, SS, ES, FLAGS,
        CF, PF, AF, ZF, TF, IF, DF, OF
    },
    sensitive=false,
    morecomment=[l]{;},
    morestring=[b]',
    morestring=[b]"
}

% ---------------------------------------------------------
% 2. Gaya kode assembly untuk slide
% ---------------------------------------------------------
\lstdefinestyle{asmslide}{
    language=MASM,
    basicstyle=\ttfamily\scriptsize,
    keywordstyle=\color{obeBlue}\bfseries,
    commentstyle=\color{obeGreen}\itshape,
    stringstyle=\color{obeMaroon},
    numberstyle=\tiny\color{obeGray},
    numbers=left,
    numbersep=5pt,
    stepnumber=1,
    backgroundcolor=\color{codeBack},
    frame=single,
    rulecolor=\color{codeFrame},
    framesep=4pt,
    breaklines=true,
    breakatwhitespace=false,
    % CATATAN PENTING: opsi `postbreak` TIDAK boleh dipakai di sini.
    % Di beamer, `breaklines=true` bersama `postbreak=<kode>` membuat
    % listings gagal dirakit dengan galat
    %   "Missing number, treated as zero <to be read again> \kern".
    % Karena \lstset{style=asmslide} bersifat global, opsi postbreak
    % juga bocor ke gaya lain (mis. outslide/keluaran) dan merusaknya.
    % Di buku ajar (kelas `book`) justru sebaliknya: postbreak aman dan
    % dipakai. Jadi larangan ini khusus untuk slide.
    showstringspaces=false,
    tabsize=4,
    columns=fullflexible,
    keepspaces=true,
    upquote=true,
    xleftmargin=12pt,
    aboveskip=0.4em,
    belowskip=0.4em
}

% ---------------------------------------------------------
% 3. Gaya keluaran program / sesi debugger untuk slide
% ---------------------------------------------------------
\lstdefinestyle{outslide}{
    basicstyle=\ttfamily\scriptsize,
    keywordstyle=\color{black},
    commentstyle=\color{obeGray},
    stringstyle=\color{black},
    numbers=none,
    backgroundcolor=\color{obeGrayLight},
    frame=single,
    rulecolor=\color{codeFrame},
    framesep=4pt,
    breaklines=true,
    showstringspaces=false,
    columns=fullflexible,
    keepspaces=true,
    upquote=true,
    xleftmargin=4pt,
    aboveskip=0.4em,
    belowskip=0.4em
}

\lstset{style=asmslide}

% ---------------------------------------------------------
% 4. Environment ringkas untuk dipakai di dalam frame
%
% Frame yang memuat kedua environment ini WAJIB memakai
% opsi [fragile], karena isinya verbatim.
% ---------------------------------------------------------
\lstnewenvironment{kodeasm}[1][]{%
    \lstset{style=asmslide, #1}%
}{}

\lstnewenvironment{keluaran}[1][]{%
    \lstset{style=outslide, #1}%
}{}


%
% Palet warna disalin dari buku_ajar/config/styles.tex agar
% identitas visual slide sama dengan buku ajar.
% =========================================================

% ---------------------------------------------------------
% 1. Paket
% ---------------------------------------------------------
\usepackage[T1]{fontenc}
\usepackage[utf8]{inputenc}
\usepackage[indonesian]{babel}
\usepackage{lmodern}
\usepackage{tikz}
\usetikzlibrary{arrows.meta, positioning, calc, shapes.geometric}
\usepackage{booktabs}
\usepackage{array}
\usepackage{amsmath}
\usepackage{amsfonts}
\usepackage{pifont}          % \ding untuk tanda checklist
\usepackage{listings}

% ---------------------------------------------------------
% 2. Palet warna OBE (identik dengan buku)
% ---------------------------------------------------------
\definecolor{obeBlue}{HTML}{1F3864}
\definecolor{obeBlueLight}{HTML}{DCE4F5}
\definecolor{obeTeal}{HTML}{0F6E6E}
\definecolor{obeTealLight}{HTML}{D9F0EE}
\definecolor{obeGreen}{HTML}{2E6B34}
\definecolor{obeGreenLight}{HTML}{DEEFDF}
\definecolor{obeAmber}{HTML}{8A5A00}
\definecolor{obeAmberLight}{HTML}{FDF1D6}
\definecolor{obeMaroon}{HTML}{7A2233}
\definecolor{obeMaroonLight}{HTML}{F7E1E5}
\definecolor{obeGray}{HTML}{4A4A4A}
\definecolor{obeGrayLight}{HTML}{F0F0F0}
\definecolor{codeBack}{HTML}{F7F7F4}
\definecolor{codeFrame}{HTML}{C9C9C0}

% ---------------------------------------------------------
% 3. Identitas mata kuliah (dipakai seluruh deck)
% ---------------------------------------------------------
\newcommand{\NamaMK}{Bahasa Pemrograman Tingkat Rendah}
\newcommand{\NamaMKShort}{Bhs.\ Pemrograman Tingkat Rendah}
\newcommand{\KodeMK}{MWCA010}
\newcommand{\Prodi}{Program Studi Teknik Informatika}
\newcommand{\Fakultas}{Fakultas Teknologi Informasi}
\newcommand{\PerguruanTinggi}{Universitas Bale Bandung}
\newcommand{\DosenPengampu}{Cecep Suwanda, S.Si., M.Kom.}
\newcommand{\SemesterTeks}{Semester Ganjil}

% \PertemuanTeks dan \JudulBab ditetapkan di masing-masing
% berkas minggu-NN.tex sebelum \begin{document}.
\providecommand{\PertemuanTeks}{Pertemuan}

% ---------------------------------------------------------
% 4. Tema beamer
% ---------------------------------------------------------
\usetheme{default}
\setbeamertemplate{navigation symbols}{}

\setbeamercolor{normal text}{fg=black!88}
\setbeamercolor{structure}{fg=obeBlue}
\setbeamercolor{alerted text}{fg=obeMaroon}

% --- judul slide: pita biru penuh lebar ---
\setbeamercolor{frametitle}{fg=white,bg=obeBlue}
\setbeamerfont{frametitle}{series=\bfseries,size=\large}
\setbeamertemplate{frametitle}{%
    \nointerlineskip%
    \begin{beamercolorbox}[%
        wd=\paperwidth,ht=2.9ex,dp=1.5ex,
        leftskip=0.4cm,rightskip=0.4cm]{frametitle}%
        \usebeamerfont{frametitle}\insertframetitle%
    \end{beamercolorbox}%
    \nointerlineskip%
    \hbox{\textcolor{obeTeal}{\rule{\paperwidth}{1.1pt}}}%
}

% --- kaki slide: identitas mata kuliah, pertemuan, nomor ---
\setbeamercolor{footline}{fg=white,bg=obeGray}
\setbeamerfont{footline}{size=\scriptsize,family=\sffamily}
\setbeamertemplate{footline}{%
    \leavevmode%
    \hbox{%
        \begin{beamercolorbox}[wd=0.42\paperwidth,ht=2.1ex,dp=1.1ex,
            leftskip=0.4cm]{footline}%
            \usebeamerfont{footline}%
            \KodeMK\ \textbullet\ \NamaMKShort%
        \end{beamercolorbox}%
        \begin{beamercolorbox}[wd=0.32\paperwidth,ht=2.1ex,dp=1.1ex,
            center]{footline}%
            \usebeamerfont{footline}\PertemuanTeks%
        \end{beamercolorbox}%
        \begin{beamercolorbox}[wd=0.26\paperwidth,ht=2.1ex,dp=1.1ex,
            rightskip=0.4cm]{footline}%
            \usebeamerfont{footline}%
            \insertframenumber\,/\,\inserttotalframenumber%
        \end{beamercolorbox}%
    }%
    \vskip0pt%
}

% --- daftar ---
\setbeamercolor{itemize item}{fg=obeTeal}
\setbeamercolor{itemize subitem}{fg=obeTeal}
\setbeamercolor{enumerate item}{fg=obeBlue}
\setbeamertemplate{itemize item}{\raise0.2ex\hbox{\tiny$\blacksquare$}}
\setbeamertemplate{itemize subitem}{\raise0.2ex\hbox{\scriptsize$\bullet$}}

% --- kotak (pengganti tcolorbox pada buku) ---
\setbeamercolor{block title}{bg=obeBlueLight,fg=obeBlue}
\setbeamercolor{block body}{bg=obeBlueLight!45,fg=black!88}
\setbeamercolor{block title alerted}{bg=obeMaroonLight,fg=obeMaroon}
\setbeamercolor{block body alerted}{bg=obeMaroonLight!45,fg=black!88}
\setbeamercolor{block title example}{bg=obeTealLight,fg=obeTeal}
\setbeamercolor{block body example}{bg=obeTealLight!45,fg=black!88}

% ---------------------------------------------------------
% 5. Makro notasi (sama persis dengan buku)
%
% Catatan: \note TIDAK didefinisikan di sini karena beamer
% sudah memakai \note untuk catatan pembicara. Untuk catatan
% di slide dipakai \catatan.
% ---------------------------------------------------------
\newcommand{\code}[1]{\texttt{#1}}
\newcommand{\asm}[1]{\texttt{#1}}
\newcommand{\reg}[1]{\texttt{#1}}
\newcommand{\flag}[1]{\texttt{#1}}
\newcommand{\dir}[1]{\texttt{#1}}
\newcommand{\file}[1]{\texttt{#1}}
\newcommand{\eng}[1]{\textit{#1}}

\newcommand{\hex}[1]{\texttt{#1h}}
\newcommand{\bin}[1]{\texttt{#1b}}
\newcommand{\dec}[1]{\texttt{#1}}

\newcommand{\important}[1]{\textbf{Penting:} #1}
\newcommand{\catatan}[1]{\textbf{Catatan:} #1}
\newcommand{\term}[1]{\textbf{#1}}

\newcommand{\CPL}[1]{\textbf{CPL#1}}
\newcommand{\CPMK}[1]{\textbf{CPMK#1}}
\newcommand{\SubCPMK}[1]{\textbf{Sub-CPMK #1}}

\newcommand{\checkbox}{\ding{113}\hspace{0.35em}}
\newcommand{\checked}{\ding{51}\hspace{0.35em}}
\newcommand{\crossed}{\ding{55}\hspace{0.35em}}

\newcommand{\flowarrow}{\ensuremath{\rightarrow}}

% ---------------------------------------------------------
% 6. Makro tata letak slide
% ---------------------------------------------------------

% Keterangan sumber materi (ketertelusuran ke buku)
\newcommand{\Sumber}[1]{%
    \vfill
    {\scriptsize\color{obeGray}Sumber: #1\par}%
}

% Frame judul deck: \JudulDeck{<judul bab>}{<pertemuan>}
\newcommand{\JudulDeck}[2]{%
\begin{frame}[plain]
    \vspace{0.25cm}
    {\color{obeBlue}\sffamily\bfseries\small \NamaMK\ (\KodeMK)\par}
    \vspace{0.12cm}
    {\color{obeGray}\sffamily\scriptsize
     \Prodi\ \textbullet\ \Fakultas\ \textbullet\ \PerguruanTinggi\par}
    \vfill
    {\color{obeBlue}\sffamily\LARGE\bfseries #1\par}
    \vspace{0.4cm}
    {\color{obeTeal}\sffamily\large #2\par}
    \vspace{0.45cm}
    {\color{obeTeal}\rule{\textwidth}{1pt}}
    \vfill
    {\color{obeGray}\sffamily\small Dosen Pengampu\par}
    \vspace{0.08cm}
    {\sffamily\bfseries \DosenPengampu\par}
    \vspace{0.2cm}
    {\color{obeGray}\sffamily\small \SemesterTeks\par}
    \vspace{0.25cm}
\end{frame}
}

% Frame pemisah bagian
\newcommand{\PemisahBagian}[1]{%
\begin{frame}[plain]
    \vfill
    \begin{center}
        {\color{obeTeal}\sffamily\small\bfseries BAGIAN\par}
        \vspace{0.35cm}
        {\color{obeBlue}\sffamily\Large\bfseries #1\par}
        \vspace{0.35cm}
        {\color{obeTeal}\rule{0.55\textwidth}{0.9pt}}
    \end{center}
    \vfill
\end{frame}
}

% ---------------------------------------------------------
% 7. Gaya penampil kode
% ---------------------------------------------------------
% =========================================================
% ppt/src/listing-ppt.tex
% Gaya penampil kode untuk slide.
%
% Diturunkan dari buku_ajar/config/listings.tex, tetapi
% disesuaikan untuk slide: ukuran lebih kecil (karena frame
% 16:9 dan kode harus utuh terlihat) dan jarak lebih rapat.
% =========================================================

% ---------------------------------------------------------
% 1. Definisi bahasa MASM/TASM 8086 (salinan dari buku)
% ---------------------------------------------------------
\lstdefinelanguage{MASM}{
    morekeywords={
        MOV, XCHG, LEA, LDS, LES, PUSH, POP, PUSHA, POPA, PUSHF, POPF,
        LAHF, SAHF, XLAT, IN, OUT, CBW, CWD,
        ADD, ADC, SUB, SBB, INC, DEC, NEG, CMP, MUL, IMUL, DIV, IDIV,
        AAA, AAS, AAM, AAD, DAA, DAS,
        AND, OR, XOR, NOT, TEST, SHL, SHR, SAL, SAR, ROL, ROR, RCL, RCR,
        JMP, JE, JZ, JNE, JNZ, JG, JGE, JL, JLE, JA, JAE, JB, JBE,
        JC, JNC, JS, JNS, JO, JNO, JP, JPE, JNP, JPO, JCXZ,
        LOOP, LOOPE, LOOPZ, LOOPNE, LOOPNZ, CALL, RET, INT, IRET, INTO,
        HLT, NOP, CLI, STI, CLD, STD, CLC, STC, CMC, WAIT, ESC, LOCK,
        DB, DW, DD, DQ, DT, EQU, ORG, SEGMENT, ENDS, ASSUME, END,
        PROC, ENDP, MACRO, ENDM, LOCAL, PUBLIC, EXTRN, EXTERN,
        INCLUDE, INCLUDELIB, MODEL, STACK, CODE, DATA, CONST, ALIGN, EVEN,
        OFFSET, PTR, BYTE, WORD, DWORD, QWORD, TBYTE,
        NEAR, FAR, SHORT, LABEL, TYPE, SIZE, LENGTH, DUP, THIS,
        AX, BX, CX, DX, AL, AH, BL, BH, CL, CH, DL, DH,
        SP, BP, SI, DI, IP, CS, DS, SS, ES, FLAGS,
        CF, PF, AF, ZF, TF, IF, DF, OF
    },
    sensitive=false,
    morecomment=[l]{;},
    morestring=[b]',
    morestring=[b]"
}

% ---------------------------------------------------------
% 2. Gaya kode assembly untuk slide
% ---------------------------------------------------------
\lstdefinestyle{asmslide}{
    language=MASM,
    basicstyle=\ttfamily\scriptsize,
    keywordstyle=\color{obeBlue}\bfseries,
    commentstyle=\color{obeGreen}\itshape,
    stringstyle=\color{obeMaroon},
    numberstyle=\tiny\color{obeGray},
    numbers=left,
    numbersep=5pt,
    stepnumber=1,
    backgroundcolor=\color{codeBack},
    frame=single,
    rulecolor=\color{codeFrame},
    framesep=4pt,
    breaklines=true,
    breakatwhitespace=false,
    % CATATAN PENTING: opsi `postbreak` TIDAK boleh dipakai di sini.
    % Di beamer, `breaklines=true` bersama `postbreak=<kode>` membuat
    % listings gagal dirakit dengan galat
    %   "Missing number, treated as zero <to be read again> \kern".
    % Karena \lstset{style=asmslide} bersifat global, opsi postbreak
    % juga bocor ke gaya lain (mis. outslide/keluaran) dan merusaknya.
    % Di buku ajar (kelas `book`) justru sebaliknya: postbreak aman dan
    % dipakai. Jadi larangan ini khusus untuk slide.
    showstringspaces=false,
    tabsize=4,
    columns=fullflexible,
    keepspaces=true,
    upquote=true,
    xleftmargin=12pt,
    aboveskip=0.4em,
    belowskip=0.4em
}

% ---------------------------------------------------------
% 3. Gaya keluaran program / sesi debugger untuk slide
% ---------------------------------------------------------
\lstdefinestyle{outslide}{
    basicstyle=\ttfamily\scriptsize,
    keywordstyle=\color{black},
    commentstyle=\color{obeGray},
    stringstyle=\color{black},
    numbers=none,
    backgroundcolor=\color{obeGrayLight},
    frame=single,
    rulecolor=\color{codeFrame},
    framesep=4pt,
    breaklines=true,
    showstringspaces=false,
    columns=fullflexible,
    keepspaces=true,
    upquote=true,
    xleftmargin=4pt,
    aboveskip=0.4em,
    belowskip=0.4em
}

\lstset{style=asmslide}

% ---------------------------------------------------------
% 4. Environment ringkas untuk dipakai di dalam frame
%
% Frame yang memuat kedua environment ini WAJIB memakai
% opsi [fragile], karena isinya verbatim.
% ---------------------------------------------------------
\lstnewenvironment{kodeasm}[1][]{%
    \lstset{style=asmslide, #1}%
}{}

\lstnewenvironment{keluaran}[1][]{%
    \lstset{style=outslide, #1}%
}{}



\renewcommand{\PertemuanTeks}{Pertemuan 1}

\begin{document}

% =========================================================
% 1. JUDUL
% =========================================================
\JudulDeck{Pengantar Assembly dan\\ Sistem Bilangan}{Pertemuan 1}

% =========================================================
% 2. CAPAIAN DAN TUJUAN
% =========================================================
\begin{frame}{Capaian dan Tujuan Pembelajaran}
\begin{block}{Sub-CPMK 1.1}
\small Menjelaskan peran bahasa assembly serta mengonversi bilangan
biner, desimal, oktal, dan heksadesimal, termasuk membedakan bilangan
bertanda dan tidak bertanda.
\end{block}

\vspace{0.1em}
{\small\textbf{CPMK yang didukung:} \CPMK{01}\quad\textbullet\quad
\textbf{CPL yang didukung:} \CPL{02}\par}

\vspace{0.15em}
\begin{exampleblock}{Setelah pertemuan ini, mahasiswa mampu:}
\footnotesize
\begin{enumerate}
    \item menjelaskan kedudukan bahasa assembly di antara bahasa mesin
          dan bahasa tingkat tinggi;
    \item menjelaskan alasan bahasa assembly masih dipakai;
    \item mengonversi bilangan biner, desimal, oktal, dan heksadesimal
          beserta langkah kerjanya;
    \item menuliskan literal bilangan yang diterima assembler;
    \item membedakan bilangan bertanda dan tidak bertanda,
          menerapkan \eng{two's complement}, serta menjelaskan
          hubungan kode ASCII dengan isi memori.
\end{enumerate}
\end{exampleblock}
\end{frame}

% =========================================================
% 3. PETA MATERI
% =========================================================
\begin{frame}{Peta Materi}
\small
\begin{enumerate}
    \item \textbf{Kedudukan dan Peran Bahasa Assembly}\\
          {\footnotesize tiga tingkat bahasa; alasan assembly masih
          dipakai; 8086 dan perangkat praktik}

    \vspace{0.25em}
    \item \textbf{Sistem Bilangan dan Konversi}\\
          {\footnotesize notasi posisional; empat basis; konversi
          antarjenis bilangan; literal assembler}

    \vspace{0.25em}
    \item \textbf{Bilangan Bertanda, Tak Bertanda, dan Data Karakter}\\
          {\footnotesize bit tanda; \eng{two's complement};
          jangkauan; kode ASCII}
\end{enumerate}

\vspace{0.35em}
\begin{block}{Praktik pada pertemuan ini}
\small Konversi bilangan dan pemeriksaan data di memori dengan
DEBUG serta emu8086.
\end{block}
\end{frame}

% =========================================================
% BAGIAN 1
% =========================================================
\PemisahBagian{Kedudukan dan Peran Bahasa Assembly}

% ---------------------------------------------------------
\begin{frame}{Tiga Tingkat Bahasa Pemrograman}
\begin{center}
\begin{tikzpicture}[
    lapis/.style={
        draw=obeBlue, line width=0.8pt, rounded corners=3pt,
        text width=10.4cm, minimum height=0.82cm, align=center,
        inner sep=5pt, font=\small
    }
]
    \node[lapis, fill=obeBlueLight] (high) {
        \textbf{Bahasa tingkat tinggi}\\[1pt]
        {\footnotesize satu baris $\Rightarrow$ banyak instruksi mesin}};
    \node[lapis, fill=obeTealLight, below=0.18cm of high] (asm) {
        \textbf{Bahasa assembly}\\[1pt]
        {\footnotesize satu instruksi $\Rightarrow$ satu instruksi mesin}};
    \node[lapis, fill=obeAmberLight, below=0.18cm of asm] (machine) {
        \textbf{Bahasa mesin}\\[1pt]
        {\footnotesize deretan bit}};
    \node[lapis, draw=obeGray, fill=obeGrayLight,
          below=0.18cm of machine] (hw) {
        \textbf{Perangkat keras}\\[1pt]
        {\footnotesize gerbang logika, bus, register}};
\end{tikzpicture}
\end{center}

\vspace{0.2em}
\begin{block}{Hubungan assembly dan bahasa mesin}
\small Bersifat \eng{one-to-one}: satu instruksi assembly
diterjemahkan menjadi tepat satu instruksi mesin.
\end{block}
\Sumber{Bab III, Bagian 1}
\end{frame}

% ---------------------------------------------------------
\begin{frame}{Mengapa Bahasa Assembly Masih Dipakai}
\small
\begin{itemize}
    \item \textbf{Kendali perangkat keras.} Akses langsung ke
          register, memori, dan port masukan/keluaran.

    \vspace{0.18em}
    \item \textbf{Ukuran dan kecepatan.} Program dapat dibuat kecil
          dan hemat instruksi pada sistem bermemori terbatas.

    \vspace{0.18em}
    \item \textbf{Pemahaman prosesor.} Instruksi, perpindahan data,
          dan pengambilan keputusan berdasarkan \eng{flag} terlihat
          langsung.

    \vspace{0.18em}
    \item \textbf{Analisis program.} Perilaku program dapat ditelaah
          pada level perangkat keras.
\end{itemize}

\vspace{0.25em}
\begin{alertblock}{Yang perlu diingat}
\small Assembly bukan pengganti bahasa tingkat tinggi; keduanya alat
dengan tujuan berbeda.
\end{alertblock}
\end{frame}

% ---------------------------------------------------------
\begin{frame}{Assembly 8086 dan Perangkat Praktik}
\small
Mata kuliah ini memakai prosesor Intel 8086 mode real 16-bit.

\vspace{0.35em}
\begin{center}
\begin{tabular}{@{}p{2.8cm}p{9.6cm}@{}}
    \toprule
    \textbf{Perangkat} & \textbf{Peran} \\
    \midrule
    TASM & Meng-\eng{assemble} berkas \file{.ASM} menjadi
    \file{.OBJ} \\
    \addlinespace
    TLINK & Mengaitkan \file{.OBJ} menjadi \file{.EXE} atau
    \file{.COM} \\
    \addlinespace
    DEBUG & Memeriksa register dan isi memori \\
    \addlinespace
    emu8086 & Menulis dan menjalankan program secara terpandu \\
    \bottomrule
\end{tabular}
\end{center}

\vspace{0.25em}
{\footnotesize\color{obeGray}Fokus bab ini: kedudukan bahasa assembly
dan cara merepresentasikan bilangan.\par}
\Sumber{Bab III, Bagian 1}
\end{frame}

% ---------------------------------------------------------
\begin{frame}{Perbandingan Tiga Tingkat}
\footnotesize
\begin{center}
\begin{tabular}{@{}>{\raggedright\arraybackslash}p{2.4cm}
    >{\raggedright\arraybackslash}p{3.2cm}
    >{\raggedright\arraybackslash}p{3.4cm}
    >{\raggedright\arraybackslash}p{3.2cm}@{}}
    \toprule
    \textbf{Aspek} & \textbf{Bahasa Mesin} &
    \textbf{Assembly} & \textbf{Tingkat Tinggi} \\
    \midrule
    Satuan penulisan & Deretan bit &
    Satu \eng{mnemonic} per instruksi & Ekspresi dan pernyataan \\
    \addlinespace
    Contoh & \code{B8 05 00} & \code{MOV AX, 5} & \code{x = 5;} \\
    \addlinespace
    Kedekatan ke mesin & Langsung & Sangat dekat & Jauh \\
    \addlinespace
    Kendali perangkat keras & Penuh & Penuh & Terbatas \\
    \addlinespace
    Portabilitas & Rendah & Rendah & Tinggi \\
    \addlinespace
    Kemudahan penulisan & Sangat sulit & Sulit & Mudah \\
    \bottomrule
\end{tabular}
\end{center}
\end{frame}

% ---------------------------------------------------------
\begin{frame}{Rangkuman Bagian 1}
\small
\begin{itemize}
    \item Tiga tingkat bahasa: mesin, assembly, dan tingkat tinggi.
          Assembly berada tepat di atas bahasa mesin dan bersifat
          \eng{one-to-one}.

    \vspace{0.15em}
    \item Semakin rendah tingkatnya, semakin dekat ke mesin dan
          semakin penuh kendali atas perangkat keras; portabilitas
          dan kemudahan penulisan menurun.

    \vspace{0.15em}
    \item Assembly dipakai untuk kendali perangkat keras, ukuran dan
          kecepatan kode, pemahaman prosesor, serta analisis program.

    \vspace{0.15em}
    \item Praktik memakai 8086 mode real 16-bit dengan TASM, TLINK,
          DEBUG, dan emu8086.

    \vspace{0.15em}
    \item Fokus bab ini adalah kedudukan assembly dan representasi
          bilangan, bukan penulisan program utuh.
\end{itemize}
\end{frame}

% =========================================================
% BAGIAN 2
% =========================================================
\PemisahBagian{Sistem Bilangan dan Konversi}

% ---------------------------------------------------------
\begin{frame}{Notasi Posisional}
\small
Nilai suatu angka ditentukan oleh digit dan posisinya. Untuk basis
$b$, angka pada posisi ke-$i$ memiliki bobot $b^i$.

\[
d_n b^n + d_{n-1}b^{n-1} + \cdots + d_1 b^1 + d_0 b^0
\]

Contoh desimal:
\[
123_{10} = 1\cdot10^2 + 2\cdot10^1 + 3\cdot10^0 = 123.
\]

\vspace{0.15em}
\begin{block}{Aturan yang sama}
\small Aturan ini berlaku untuk semua basis; yang berubah hanya nilai
basis $b$.
\end{block}
\Sumber{Bab III, Bagian 2}
\end{frame}

% ---------------------------------------------------------
\begin{frame}{Empat Basis yang Dipakai}
\small
\begin{center}
\begin{tabular}{@{}p{3.0cm}p{4.0cm}p{4.0cm}@{}}
    \toprule
    \textbf{Basis} & \textbf{Digit} & \textbf{Contoh} \\
    \midrule
    Biner ($b=2$) & 0 dan 1 & \bin{11110011} \\
    \addlinespace
    Oktal ($b=8$) & 0 sampai 7 & $363_8$ \\
    \addlinespace
    Desimal ($b=10$) & 0 sampai 9 & 243 \\
    \addlinespace
    Heksadesimal ($b=16$) & 0--9 dan A--F & \hex{F3} \\
    \bottomrule
\end{tabular}
\end{center}

\vspace{0.3em}
{\footnotesize\color{obeGray}Pada heksadesimal: A = 10, B = 11,
C = 12, D = 13, E = 14, dan F = 15.\par}
\Sumber{Bab III, Bagian 2}
\end{frame}

% ---------------------------------------------------------
\begin{frame}{Biner ke Desimal}
\small
Kalikan setiap bit dengan bobot posisinya, lalu jumlahkan.

\[
\begin{aligned}
&1\cdot2^7 + 1\cdot2^6 + 1\cdot2^5 + 1\cdot2^4\\
&\quad + 0\cdot2^3 + 0\cdot2^2 + 1\cdot2^1 + 1\cdot2^0 = 243.
\end{aligned}
\]

\[
128 + 64 + 32 + 16 + 2 + 1 = 243.
\]

\vspace{0.15em}
\begin{block}{Hasil}
\small \bin{11110011} sama dengan 243 desimal.
\end{block}
\Sumber{Bab III, Bagian 2}
\end{frame}

% ---------------------------------------------------------
\begin{frame}{Desimal ke Biner}
\small
Bagi dengan 2 secara berulang, catat sisanya, lalu baca sisa dari
bawah ke atas.

\vspace{0.2em}
\begin{center}
\begin{tabular}{@{}p{4.2cm}p{3.2cm}p{2.2cm}@{}}
    \toprule
    \textbf{Pembagian} & \textbf{Hasil} & \textbf{Sisa} \\
    \midrule
    $243 \div 2$ & 121 & 1 \\
    $121 \div 2$ & 60 & 1 \\
    $60 \div 2$ & 30 & 0 \\
    $30 \div 2$ & 15 & 0 \\
    $15 \div 2$ & 7 & 1 \\
    $7 \div 2$ & 3 & 1 \\
    $3 \div 2$ & 1 & 1 \\
    $1 \div 2$ & 0 & 1 \\
    \bottomrule
\end{tabular}
\end{center}

\vspace{0.1em}
{\small Hasil: \bin{11110011}.\par}
\Sumber{Bab III, Bagian 2}
\end{frame}

% ---------------------------------------------------------
\begin{frame}{Oktal dan Heksadesimal}
\small
Kelompokkan bit dari kanan: tiga bit untuk oktal dan empat bit untuk
heksadesimal.

\vspace{0.25em}
\begin{columns}[t]
\begin{column}{0.48\textwidth}
    {\color{obeTeal}\bfseries Oktal}\\[0.2em]
    \code{011 110 011} $\rightarrow$ 3, 6, 3\\
    Hasil: $363_8$
\end{column}
\begin{column}{0.48\textwidth}
    {\color{obeTeal}\bfseries Heksadesimal}\\[0.2em]
    \code{1111 0011} $\rightarrow$ F, 3\\
    Hasil: \hex{F3}
\end{column}
\end{columns}

\vspace{0.35em}
Pemeriksaan dengan bobot posisi:
\[
3\cdot8^2 + 6\cdot8 + 3 = 243,
\qquad
15\cdot16 + 3 = 243.
\]

\begin{block}{Ringkas}
\small Pengelompokan bit lebih cepat; bobot posisi dipakai untuk
memeriksa hasil.
\end{block}
\end{frame}

% ---------------------------------------------------------
\begin{frame}{Literal Bilangan pada Assembler}
\small
\begin{itemize}
    \item \textbf{Desimal:} \code{243} atau \code{243D}.
    \item \textbf{Biner:} \code{11110011B}.
    \item \textbf{Heksadesimal:} \code{0F3H}; bila digit pertama
          A--F, tambahkan nol di depan.
    \item \textbf{Karakter:} \code{'A'} atau \code{"A"}.
\end{itemize}

\vspace{0.3em}
\begin{alertblock}{Kesalahan yang sering terjadi}
\small Menulis \code{F3H} tanpa nol dapat dibaca assembler sebagai
nama label, bukan bilangan.
\end{alertblock}
\Sumber{Bab III, Bagian 2; [2] BAB V}
\end{frame}

% ---------------------------------------------------------
\begin{frame}{Rangkuman Bagian 2}
\small
\begin{itemize}
    \item Notasi posisional: nilai angka adalah jumlah digit
          dikalikan bobot basisnya.

    \vspace{0.15em}
    \item Empat basis yang dipakai: biner, oktal, desimal, dan
          heksadesimal.

    \vspace{0.15em}
    \item Biner ke desimal dihitung dengan bobot posisi; desimal ke
          biner dengan pembagian berulang dan sisa dibaca terbalik.

    \vspace{0.15em}
    \item Oktal dibentuk dari kelompok tiga bit; heksadesimal dari
          kelompok empat bit.

    \vspace{0.15em}
    \item Literal memakai akhiran D, B, atau H; heksadesimal yang
          diawali huruf perlu nol di depan.
\end{itemize}
\end{frame}

% =========================================================
% BAGIAN 3
% =========================================================
\PemisahBagian{Bilangan Bertanda, Tak Bertanda, dan Data Karakter}

% ---------------------------------------------------------
\begin{frame}{Bilangan Tidak Bertanda}
\small
Seluruh bit dipakai sebagai besaran. Untuk $n$ bit, jangkauannya
adalah 0 sampai $2^n - 1$.

\vspace{0.3em}
\begin{center}
\begin{tabular}{@{}p{3.0cm}p{7.0cm}@{}}
    \toprule
    \textbf{Panjang} & \textbf{Jangkauan} \\
    \midrule
    8 bit & 0 sampai 255 \\
    \addlinespace
    16 bit & 0 sampai 65.535 \\
    \bottomrule
\end{tabular}
\end{center}

\vspace{0.3em}
Contoh: \hex{FF} pada 8 bit bernilai 255; \hex{FFFF} pada 16 bit
bernilai 65.535.
\Sumber{Bab III, Bagian 3}
\end{frame}

% ---------------------------------------------------------
\begin{frame}{Bilangan Bertanda dan Two's Complement}
\small
Bit paling kiri adalah \term{bit tanda}: 0 untuk positif atau nol,
1 untuk negatif.

\vspace{0.25em}
\begin{center}
\begin{tabular}{@{}p{6.0cm}p{5.0cm}@{}}
    \toprule
    \textbf{Langkah} & \textbf{Hasil 8 bit} \\
    \midrule
    Bilangan positif 37 & \bin{00100101} \\
    \addlinespace
    Balikkan seluruh bit & \bin{11011010} \\
    \addlinespace
    Tambahkan 1 & \bin{11011011} = \hex{DB} \\
    \bottomrule
\end{tabular}
\end{center}

\vspace{0.2em}
{\footnotesize\color{obeGray}Urutan: tulis bilangan positif, balik
bit, lalu tambahkan 1.\par}
\Sumber{Bab III, Bagian 3}
\end{frame}

% ---------------------------------------------------------
\begin{frame}{Jangkauan Bilangan Bertanda}
\small
\begin{center}
\begin{tabular}{@{}p{2.6cm}p{4.6cm}p{5.2cm}@{}}
    \toprule
    \textbf{Panjang} & \textbf{Tidak Bertanda} &
    \textbf{Bertanda} \\
    \midrule
    8 bit & 0 sampai 255 & $-128$ sampai 127 \\
    \addlinespace
    16 bit & 0 sampai 65.535 & $-32.768$ sampai 32.767 \\
    \bottomrule
\end{tabular}
\end{center}

\vspace{0.3em}
Nilai \bin{10000000} pada 8 bit bertanda adalah $-128$ dan tidak
memiliki pasangan positif pada panjang bit yang sama.
\Sumber{Bab III, Bagian 3}
\end{frame}

% ---------------------------------------------------------
\begin{frame}{Menafsirkan Pola Bit}
\small
Pola bit yang sama dapat mempunyai dua makna.

\vspace{0.25em}
\begin{itemize}
    \item \hex{F3} = \bin{11110011}.
    \item Tidak bertanda: 243.
    \item Bertanda 8 bit: $243 - 256 = -13$.
\end{itemize}

\vspace{0.25em}
\begin{itemize}
    \item \eng{Sign Flag} mencatat tanda hasil operasi.
    \item Program menentukan tafsir melalui jenis operasi dan
          lompatan bersyarat yang dipilih.
\end{itemize}

\vspace{0.2em}
\begin{block}{Aturan praktis}
\small Bit tanda 0: nilai bertanda sama dengan nilai tidak bertanda.
Bit tanda 1: kurangkan $2^n$ dari nilai tidak bertanda.
\end{block}
\end{frame}

% ---------------------------------------------------------
\begin{frame}{Kode ASCII dan Data Karakter}
\small
Karakter disimpan sebagai kode numerik satu byte.

\vspace{0.3em}
\begin{center}
\begin{tabular}{@{}p{3.0cm}p{5.0cm}@{}}
    \toprule
    \textbf{Karakter} & \textbf{Kode} \\
    \midrule
    \code{'A'} & 65 = \hex{41} \\
    \addlinespace
    \code{'a'} & 97 = \hex{61} \\
    \addlinespace
    \code{'0'} & 48 = \hex{30} \\
    \addlinespace
    spasi & 32 = \hex{20} \\
    \bottomrule
\end{tabular}
\end{center}

\vspace{0.25em}
\begin{block}{String}
\small String adalah byte-byte ASCII berurutan. Menulis \code{'A'}
pada segmen data sama dengan menyimpan \hex{41}.
\end{block}
\Sumber{Bab III, Bagian 3}
\end{frame}

% ---------------------------------------------------------
\begin{frame}{Rangkuman Bagian 3}
\small
\begin{itemize}
    \item Tidak bertanda: seluruh bit sebagai besaran; 8 bit
          0--255, 16 bit 0--65.535.

    \vspace{0.15em}
    \item Bertanda: bit kiri adalah tanda; \eng{two's complement}
          dibentuk dengan membalik bit lalu menambah 1.

    \vspace{0.15em}
    \item Jangkauan bertanda: 8 bit $-128$ sampai 127; 16 bit
          $-32.768$ sampai 32.767.

    \vspace{0.15em}
    \item \hex{F3} dapat berarti 243 atau $-13$; tafsirnya ditentukan
          program dan dicatat \eng{Sign Flag}.

    \vspace{0.15em}
    \item ASCII: \code{'A'} = \hex{41}, \code{'a'} = \hex{61},
          \code{'0'} = \hex{30}, spasi = \hex{20}.
\end{itemize}
\end{frame}

% =========================================================
% CONTOH PROGRAM
% =========================================================
\begin{frame}{Contoh 1: Konversi Nilai 200}
\small
\begin{columns}[t]
\begin{column}{0.46\textwidth}
\footnotesize
\begin{tabular}{@{}p{2.4cm}p{1.8cm}p{1.0cm}@{}}
    \toprule
    \textbf{Pembagian} & \textbf{Hasil} & \textbf{Sisa} \\
    \midrule
    $200 \div 2$ & 100 & 0 \\
    $100 \div 2$ & 50 & 0 \\
    $50 \div 2$ & 25 & 0 \\
    $25 \div 2$ & 12 & 1 \\
    $12 \div 2$ & 6 & 0 \\
    $6 \div 2$ & 3 & 0 \\
    $3 \div 2$ & 1 & 1 \\
    $1 \div 2$ & 0 & 1 \\
    \bottomrule
\end{tabular}
\end{column}
\begin{column}{0.50\textwidth}
\footnotesize
Biner: \bin{11001000}\\[0.3em]
Oktal: \code{011 001 000} $\rightarrow$ $310_8$\\[0.3em]
Heksadesimal: \code{1100 1000} $\rightarrow$ \hex{C8}\\[0.3em]
Bertanda 8 bit: $200 - 256 = -56$.
\end{column}
\end{columns}

\vspace{0.35em}
\begin{block}{Satu pola, dua tafsir}
\small \bin{11001000} berarti 200 sebagai bilangan tidak bertanda atau
$-56$ sebagai bilangan bertanda.
\end{block}
\Sumber{Bab III, Contoh 1}
\end{frame}

% ---------------------------------------------------------
\begin{frame}[fragile]{Contoh 2: Menyimpan Bilangan dalam Basis Berbeda}
\begin{columns}[t]
\begin{column}{0.59\textwidth}
\begin{kodeasm}
.MODEL SMALL
.STACK 100h
.DATA
    desimal  DB 243
    biner    DB 11110011b
    heksa    DB 0F3h
    karakter DB 'A'
.CODE
START:
    MOV AX, @DATA
    MOV DS, AX
    MOV AH, 4Ch
    MOV AL, 00h
    INT 21h
END START
\end{kodeasm}
\end{column}
\begin{column}{0.38\textwidth}
{\color{obeTeal}\bfseries Isi memori}\\[0.3em]
\begin{keluaran}
C:\>debug bilangan.exe
-d ds:0
0B2A:0000  F3 F3 F3 41
\end{keluaran}

\vspace{0.4em}
{\footnotesize Tiga literal pertama menghasilkan \hex{F3}; karakter
\code{'A'} menghasilkan \hex{41}.\par}
\end{column}
\end{columns}

\vspace{0.15em}
{\footnotesize\color{obeGray}Tidak ada keluaran ke layar; data
diperiksa di memori.\par}
\Sumber{Bab III, Contoh 2}
\end{frame}

% =========================================================
% PRAKTIKUM
% =========================================================
\begin{frame}{Praktikum: Konversi Bilangan dan Pemeriksaan Data}
\begin{block}{Judul praktikum}
\small Konversi bilangan dan pemeriksaan data di memori\\[0.15em]
\textbf{Sub-CPMK:} 1.1 \quad\textbullet\quad
\textbf{Waktu:} 100 menit \quad\textbullet\quad
\textbf{Bentuk:} individu
\end{block}

\vspace{0.2em}
\begin{columns}[t]
\begin{column}{0.46\textwidth}
{\color{obeBlue}\bfseries Alat dan bahan}
\begin{itemize}\footnotesize
    \item Komputer dengan Windows
    \item emu8086
    \item DOSBox dengan TASM, TLINK, dan DEBUG
    \item Berkas \file{bilangan.asm}
\end{itemize}
\end{column}
\begin{column}{0.50\textwidth}
{\color{obeBlue}\bfseries Langkah kerja}
\begin{enumerate}\footnotesize
    \item Lengkapi tabel konversi bilangan.
    \item Tulis dan rakit \file{bilangan.asm}.
    \item Jalankan \file{debug bilangan.exe}, lalu \file{d ds:0}.
    \item Ubah data menjadi 37, \code{0DBh}, dan \code{'Z'}.
    \item Bandingkan hasil DEBUG dengan emu8086.
\end{enumerate}
\end{column}
\end{columns}

\vspace{0.2em}
{\footnotesize\textbf{Dikumpulkan:} tabel konversi, berkas
\file{.asm} modifikasi, tangkapan layar DEBUG dan emu8086, serta
analisis singkat.\par}
\Sumber{Bab III, Praktikum}
\end{frame}

% =========================================================
% LATIHAN
% =========================================================
\begin{frame}{Latihan}
\small
\begin{enumerate}
    \item Jelaskan perbedaan bahasa mesin, bahasa assembly, dan bahasa
          tingkat tinggi. Sertakan satu operasi sederhana yang ditulis
          pada ketiga tingkat.

    \vspace{0.18em}
    \item Mengapa hubungan bahasa assembly dan bahasa mesin disebut
          \eng{one-to-one}? Bandingkan dengan bahasa tingkat tinggi.

    \vspace{0.18em}
    \item Ubah 156 desimal menjadi biner, oktal, dan heksadesimal.
          Tuliskan langkah yang Anda pakai.

    \vspace{0.18em}
    \item Diberikan \bin{10110110}. Tentukan nilainya sebagai
          bilangan tidak bertanda dan sebagai bilangan bertanda 8 bit.

    \vspace{0.18em}
    \item Tentukan bentuk 8 bit dari $-45$ dengan
          \eng{two's complement}, lalu periksa hasilnya dengan
          mengurangkan 256 dari nilai tidak bertandanya.
\end{enumerate}
\end{frame}

% =========================================================
% RANGKUMAN
% =========================================================
\begin{frame}{Rangkuman}
\small
\begin{itemize}
    \item Tiga tingkat bahasa: mesin, assembly, dan tingkat tinggi.
          Assembly bersifat \eng{one-to-one} terhadap bahasa mesin.

    \vspace{0.15em}
    \item Assembly dipakai untuk kendali perangkat keras, ukuran dan
          kecepatan kode, serta pemahaman prosesor.

    \vspace{0.15em}
    \item Empat basis: biner, oktal, desimal, dan heksadesimal,
          masing-masing dengan aturan digit dan bobot posisi.

    \vspace{0.15em}
    \item Konversi memakai bobot posisi, pembagian berulang, atau
          pengelompokan tiga dan empat bit.

    \vspace{0.15em}
    \item Literal assembler memakai akhiran D, B, atau H; heksadesimal
          yang diawali huruf perlu nol di depan.

    \vspace{0.15em}
    \item \hex{F3} berarti 243 atau $-13$; \code{'A'} disimpan sebagai
          \hex{41}.
\end{itemize}
\end{frame}

% ---------------------------------------------------------
\begin{frame}{Cek Pemahaman}
\footnotesize
\begin{enumerate}
    \item Bilangan biner adalah bilangan yang berbasis \dots
          \begin{enumerate}
              \item 2 dengan digit 0 dan 1
              \item 8 dengan digit 0 sampai 7
              \item 10 dengan digit 0 sampai 9
              \item 16 dengan digit 0 sampai F
          \end{enumerate}

    \vspace{0.2em}
    \item Nilai desimal dari \hex{F3} adalah \dots
          \begin{enumerate}
              \item 153 \quad \item 243 \quad
              \item 253 \quad \item 263
          \end{enumerate}

    \vspace{0.2em}
    \item Pola 8 bit \bin{11111111} sebagai bilangan bertanda
          bernilai \dots
          \begin{enumerate}
              \item 255 \quad \item 127 \quad
              \item $-1$ \quad \item $-128$
          \end{enumerate}
\end{enumerate}
\end{frame}

% =========================================================
% PENUTUP
% =========================================================
\begin{frame}{Penutup}
\begin{block}{Sebelum pertemuan berikutnya}
\small
\begin{itemize}
    \item Selesaikan praktikum: tabel konversi, \file{bilangan.asm}
          modifikasi, tangkapan layar, dan analisis.
    \item Baca \textbf{Bab IV} — Arsitektur Mikroprosesor 8086:
          BIU dan EU.
    \item Siapkan pertanyaan tentang register dan alur pengambilan
          instruksi.
\end{itemize}
\end{block}

\vspace{0.4em}
\begin{center}
{\color{obeBlue}\sffamily\large\bfseries Terima kasih\par}
\vspace{0.15em}
{\color{obeGray}\sffamily\small \NamaMK\ (\KodeMK)
\textbullet\ \SemesterTeks\par}
\end{center}
\end{frame}

\end{document}
